\documentclass[a4paper]{article}
\usepackage{amsfonts}
\usepackage{amsmath}
\usepackage{amssymb}
\usepackage{graphicx}

\newcommand{\dint}{\displaystyle\int}

\begin{document}
  
\noindent \large Auteur: Michael Livchits \\
\noindent \large Studierichting: Informatica\\
\noindent \large Oefeningenreeks 5F nr. 1.6\\

\medskip

\normalsize

$r = -4\sin 3\theta$\\

$r' = -12\cos 3\theta$\\

\begin{tabular}{c|ccccccc}
$\theta$ & $0$ & $\dfrac{\pi}{6}$ & $\dfrac{\pi}{3}$ & $\dfrac{\pi}{2}$ & $\dfrac{2\pi}{3}$ & $\dfrac{5\pi}{6}$ & $\pi$\\ \hline
$r$ & $0$ & $-$ & $0$ & $+$ & $0$ & $-$ & $0$ \\
$r'$ & $-$ & $0$ & $+$ & $0$ & $-$& $0$ & $+$ \\ \hline
$r$ & $\searrow$ &$m(-4)$& $\nearrow$ & $M(4)$ & $\searrow$&$m(-4)$& $\nearrow$ \\
\end{tabular}\\

\begin{figure}[h]
\includegraphics[width=6cm]{5F-1.6-michael.livchits.INF.png} \\
\end{figure}

$S_1 = \dfrac{1}{2} \displaystyle\int_{0}^{\tfrac{\pi}{3}} (-4\sin 3\theta)^2 d\theta$\\

$= \dfrac{1}{2} \displaystyle\int_{0}^{\tfrac{\pi}{3}} 16\sin^2 3\theta d\theta$\\

$= 8\displaystyle\int_{0}^{\tfrac{\pi}{3}} \sin^2 3\theta d\theta$\\

$= 8\displaystyle\int_{0}^{\tfrac{\pi}{3}} \dfrac{1-\cos6\theta}{2}d\theta$\\

$= \left[4\theta - \dfrac{2}{3}\sin 6\theta\right]_{0}^{\tfrac{\pi}{3}}$\\

$= \dfrac{4\pi}{3}$\\

$S = S_1\cdot3$\\

$= 4\pi$

\begin{thebibliography}{999}
%\bibitem{a} Referentie
\end{thebibliography}
\end{document} 